\documentclass[10pt, ngerman, a4paper,toc=graduated, parskip=half, numbers=noenddot, headings=standardclasses, nonumberlist]{scrartcl}
\usepackage[utf8]{inputenc}
\usepackage[T1]{fontenc}
\usepackage{tgheros}
\renewcommand{\familydefault}{\sfdefault}
\usepackage{babel}
%\usepackage{lmodern}
\usepackage[top=2.5cm, right=2.5cm, bottom=2cm, left=2.5cm, footskip=0.5cm]{geometry}
\usepackage{calc}
\usepackage{graphicx}
\usepackage{svg}
\usepackage{amsmath}
\usepackage{tabularx}
\usepackage{booktabs}
\usepackage[official]{eurosym}
\usepackage{lastpage}
\usepackage{etoolbox}
\usepackage{amssymb}
\usepackage{enumitem}
%\usepackage{natbib}
%\bibliographystyle{agsm}
\bibliographystyle{IEEETran}
\usepackage[hidelinks]{hyperref}
\usepackage{dramatist}
\usepackage{listings}

\usepackage[table]{xcolor}
\usepackage{multirow}
\usepackage{array}

\usepackage[dvipsnames]{xcolor}
\definecolor{beispielfarbe}{RGB}{28, 54, 97}

\usepackage{adjustbox,lipsum}
\usepackage{makecell}
\usepackage{caption}
\usepackage{longtable}
\usepackage{tabularx}
\usepackage{ragged2e}
\usepackage{array}
\newcolumntype{Y}{>{\RaggedRight\arraybackslash}X}

\lstset{ 
    breakatwhitespace=false,         % sets if automatic breaks should only happen at whitespace
    breaklines=true,                 % sets automatic line breaking
    frame=single,                    % adds a frame around the code
    keepspaces=true,                 % keeps spaces in text, useful for keeping indentation of code (possibly needs columns=flexible)
    numbers=left,                    % where to put the line-numbers; possible values are (none, left, right)
    numbersep=5pt,                   % how far the line-numbers are from the code
    numberstyle=\color{gray},        % the style that is used for the line-numbers
    rulecolor=\color{black}         % if not set, the frame-color may be changed on line-breaks within not-black text (e.g. comments (green here))
}

\lstdefinelanguage{json}{
    basicstyle=\ttfamily\small,
    numbers=left,
    numberstyle=\color{gray},
    stepnumber=1,
    numbersep=8pt,
    showstringspaces=false,
    breaklines=true,
    frame=single,
    literate=
     *{0}{{{\color{black}0}}}{1}
      {1}{{{\color{black}1}}}{1}
      {2}{{{\color{black}2}}}{1}
      {3}{{{\color{black}3}}}{1}
      {4}{{{\color{black}4}}}{1}
      {5}{{{\color{black}5}}}{1}
      {6}{{{\color{black}6}}}{1}
      {7}{{{\color{black}7}}}{1}
      {8}{{{\color{black}8}}}{1}
      {9}{{{\color{black}9}}}{1}
      {:}{{{\color{black}{:}}}}{1}
      {,}{{{\color{black}{,}}}}{1}
      {\{}{{{\color{black}{\{}}}}{1}
      {\}}{{{\color{black}{\}}}}}{1}
      {[}{{{\color{black}{[}}}}{1}
      {]}{{{\color{black}{]}}}}{1},
}

\lstdefinelanguage{JavaScript}{
    keywords={function, return, const, let, var, if, else, for, while, new, true, false, null},
    sensitive=true,
    comment=[l]{//},
    morecomment=[s]{/*}{*/},
    morestring=[b]",
    morestring=[b]',
}

\usepackage[acronym, toc]{glossaries}
\makenoidxglossaries


% ------------- Beispiele für Abkürzungen
\newacronym{erd}{ERD}{Entity-Relation-Diagram}
\newacronym{mvp}{MVP}{Minimum Viable Product}


% ------------- Beispiele für Glossareinträge
\newglossaryentry{gitlabrunner}{name={GitLab-Runner},description={Anwendung, die auf Aktionen eines Git Repositories reagiert und vorkonfigurierte Jobs ausführt}}
\newglossaryentry{rest}{name={REST}, description = {Representational State Transfer, Ein Paradigma für die Architektur von \gls{crud} Webanwendungen }}


% ------------- Abbildungen mit Tikz erstellen
\usepackage{tikz}
\usetikzlibrary{arrows,automata,shapes,calc}


% ------------- Zeilenabstand
\usepackage{setspace}
\singlespacing

% ------------- Kopfzeile
\usepackage[automark]{scrlayer-scrpage}
\pagestyle{scrheadings}
\clearscrheadings
\ihead{}
%\ihead[\rightmark]{\pagemark}
\chead{}
%\chead{}{}
%\ohead{\pagemark}
%\ohead[\pagemark]{\rightmark}
\ofoot{\pagemark}
%\ofoot[\pagemark]{\rightmark}


\renewcommand{\contentsname}{Inhaltsverzeichnis}
\renewcommand{\listfigurename}{Abbildungsverzeichnis}
\renewcommand{\listtablename}{Tabellenverzeichnis}

\RedeclareSectionCommand[
  font=\normalfont\bfseries
]{section}

\RedeclareSectionCommand[
  font=\normalfont\bfseries
]{subsection}

\RedeclareSectionCommand[
  font=\normalfont\bfseries
]{subsubsection}

%-------------- Ende allgemeine Angaben -----------------------------
\begin{document}

%-------------- Beginn Titelseite -----------------------------------
\hypersetup{pageanchor=false}
\thispagestyle{empty}

\vspace{1cm}

\begin{figure}[htbp]
    \centering

    \begin{minipage}[t]{0.45\textwidth}
        \raggedright
        \includesvg[width=5cm,height=10cm,keepaspectratio]{images/logo_name}
        \captionsetup{justification=raggedright,singlelinecheck=false}
        \caption{Logo von xxx \cite{citation_name}}
        \label{fig:abbildung_label}
    \end{minipage}
    \hfill
    \begin{minipage}[t]{0.45\textwidth}
        \raggedleft
        \includesvg[width=5cm,height=10cm,keepaspectratio]{images/logo2_name}
        \captionsetup{justification=raggedleft,singlelinecheck=false}
        \caption{Logo von yyy \cite{citation_name}}
        \label{fig:abbildung2_label}
    \end{minipage}

\end{figure}

\vspace{1cm}
\begin{center}
	{\Large Art des Arbeit} \\
	\vspace{3cm}
	{\huge \textbf{Titel}}\\
    \vspace{0.3cm}
    {\LARGE \textbf{{Untertitel}}}\\
    
\end{center}

\vspace{1.5cm}
\begin{center}
	Abgabe: Ort, den TT.MM.JJJJ\\[4\baselineskip]
	\textbf{\textcolor{beispielfarbe}{Vorname Nachname}}

    Straße \& HNR \\
    PLZ Ort \\
    Prüflingsnummer: XXX \\
    Ausbildungsgang \\[4\baselineskip]

    Firmenname \\
    Firmenadresse \\
    Firmen-PLZ Firmen-Ort
	
	Betriebliche:r Ausbilder:in \\
    Vorname Nachname

    Fachliche:r Betreuer:in \\
    Vorname Nachname \\
\end{center}
%-------------- Ende Titelseite -------------------------------------

%-------------- Inhaltsverzeichnis
\newpage
\hypersetup{pageanchor=true}
\pagenumbering{arabic}
\renewcommand{\thepage}{\Roman{page}}
\tableofcontents{}


%-------------- Selbstständigkeitserklärung
\newpage
\cleardoublepage
\addsec{Selbstständigkeitserklärung}

Hiermit erkläre ich, dass ich die vorliegende Arbeit selbstständig und ohne fremde Hilfe angefertigt habe. 
Alle Stellen, die wörtlich oder sinngemäß aus veröffentlichten oder nicht veröffentlichten Quellen übernommen wurden, 
sind als solche kenntlich gemacht.

Ich versichere außerdem, dass die Arbeit in gleicher oder ähnlicher Form noch keiner anderen Prüfungsstelle vorgelegt wurde.

\vspace{2cm}

\noindent
\begin{tabular}{@{}p{6cm}p{6cm}@{}}
Ort, Datum & Unterschrift \\[1.5cm]
\hrulefill & \hrulefill
\end{tabular}

%-------------- Abbildungsverzeichnis
\newpage
\cleardoublepage\phantomsection
\addcontentsline{toc}{section}{Abbildungsverzeichnis}
\listoffigures
\renewcommand{\figurename}{Abb.}

%-------------- Abkürzungsverzeichnis
\newpage
\cleardoublepage\phantomsection
\printnoidxglossary[type=\acronymtype, title= Abkürzungsverzeichnis]

%-------------- Glossar
\newpage
\cleardoublepage\phantomsection
\printnoidxglossary[title= Glossar]

%-------------- Tabellenverzeichnis
\newpage
\cleardoublepage\phantomsection
\addcontentsline{toc}{section}{Tabellenverzeichnis}
\listoftables

\newpage
\pagenumbering{arabic}
\renewcommand{\thepage}{\arabic{page}}

%-------------- Beispiel einer Section
\section{Beispiel einer Section}
\label{beispiel_einer_section}
Beispiel einer Abkürzung zu \gls{mvp}.
%-------------- Einleitung

%-------------- Beispiel einer Subsection
\subsection{Beispiel einer Subsection}
\label{beispiel_einer_subsection}
Beispiel eines Glossareintrags \gls{gitlabrunner}
%-------------- Beispiel einer Subsection

%-------------- Beispiel einer Subsubsection
\subsection{Beispiel einer Subsubsection}
\label{beispiel_einer_subsubsection}
%-------------- Beispiel einer Subsubsection

%-------------- Literaturverzeichnis
\newpage
\addcontentsline{toc}{section}{Literatur}
\bibliography{Quellen}

%-------------- Anhang
\newpage
\appendix
\section{Anhang}

\newpage
\subsection{Beispielcode zur Normalisierung und Fuzzy-Bewertung}
\label{a1_normalisierung_fuzzy}

\begin{lstlisting}[language=JavaScript, basicstyle=\ttfamily\small, numbers=none, caption={Normalisierung und Fuzzy-Bewertung}, label={lst:normalisierung_fuzzy}]
function normalizeText(input) {
    return String(input ?? "")
        .toLowerCase()
        .normalize("NFD")
        .replace(/[\u0300-\u036f]/g, "")
		.replace(/\u00DF/g, "ss")
        .replace(/[^a-z0-9]+/g, " ")
        .trim()
        .replace(/\s+/g, " ");
}

function tokenize(norm) {
    if (!norm) return [];
    return norm.split(" ").filter(Boolean);
}

function levenshtein(a, b) {
    if (a === b) return 0;
    if (!a) return b.length;
    if (!b) return a.length;

    const v0 = new Array(b.length + 1);
    const v1 = new Array(b.length + 1);

    for (let i = 0; i <= b.length; i++) v0[i] = i;

    for (let i = 0; i < a.length; i++) {
        v1[0] = i + 1;
        for (let j = 0; j < b.length; j++) {
            const cost = a[i] === b[j] ? 0 : 1;
            v1[j + 1] = Math.min(
                v1[j] + 1,
                v0[j + 1] + 1,
                v0[j] + cost
            );
        }
        for (let j = 0; j <= b.length; j++) v0[j] = v1[j];
    }

    return v0[b.length];
}

function tokenJaccard(aTokens, bTokens) {
    if (!aTokens.length || !bTokens.length) return 0;
    const aSet = new Set(aTokens);
    const bSet = new Set(bTokens);
    let inter = 0;

    for (const t of aSet) {
        if (bSet.has(t)) inter++;
    }

    const union = new Set([...aSet, ...bSet]).size;
    return union ? inter / union : 0;
}

function scoreCandidate(queryNorm, candNorm) {
    if (!candNorm) return 0;
    if (queryNorm === candNorm) return 1;

    const qTokens = tokenize(queryNorm);
    const cTokens = tokenize(candNorm);

    let substringScore = 0;
    if (queryNorm.length >= 3) {
        if (queryNorm.includes(candNorm)) substringScore = 0.97;
        else if (candNorm.includes(queryNorm)) substringScore = 0.93;
    }

    const tokenScore = tokenJaccard(qTokens, cTokens);
    const dist = levenshtein(queryNorm, candNorm);
    const levScore = 1 - dist / Math.max(queryNorm.length, candNorm.length);

    return Math.max(
        substringScore,
        tokenScore * 0.9,
        levScore * 0.88 + tokenScore * 0.12
    );
}
\end{lstlisting}

\end{document}